\section{Invariant Distributions}

Our question in this chapter is how does the probabilities \(\Prob\pqty{X_n = x}\) or \(\Prob_x \pqty{X_n = y}\) behave as \(n \to \infty\)?

Say \(\Prob\pqty{X_n = j} \to \pi_j\) as \(n \to \infty\), for some probability distribution \(\pqty{\pi_j}_{j \in I}\). Then we expect that if \(X_n \sim \pi\), then so does \(X_{n + 1} \sim \pi\), i.e.,
\[
    \pi_j = \Prob\pqty{X_{n + 1} = j} = \sum_{i \in I} \Prob\pqty{X_n = i} \Prob\pqty{\Conditional{X_{n + 1} = j}{X_n = i}} = \sum_{i \in I} \pi_i P_{ij},
\]
i.e., taking \(\pi\) as a row vector,
\[
    \pi = \pi P
\]

\begin{definition}[Measure, Probability Distribution, Invariant Distribution]
    \begin{itemize}
        \item A \term{measure} \(\lambda\) on \(I\) is a vector \(\lambda = \pqty{\lambda_i}_{i \in I}\) with non-negative entries, i.e., \(\lambda_i \ge 0\) for all \(i \in I\).
        \item A measure \(\lambda\) is a \term{probability distribution} if \(\sum_{i \in I} \lambda_i = 1\).
        \item A measure \(\lambda\) is \term{invariant} or \term{stationary} with respect to a transition matrix \(P\) if \(\lambda P = \lambda\).
    \end{itemize}
\end{definition}

\begin{theorem}[Invariance of Distribution]
    If \(\pi\) is invariant for \(X\), and \(X_0 \sim \pi\), then \(X_n \sim \pi\) for all \(n \geq 0\).
\end{theorem}

\begin{proof}
    We show by induction. This is obvious for \(n = 0\). Now, suppose \(X_N \sim \pi\), then
    \begin{align*}
        \Prob\pqty{X_{n + 1} = j} &= \sum_{i \in I} \Prob\pqty{X_n = i, X_{n + 1} = j} \\
        &= \sum_{i \in I} \pi_i P_{ij} \\
        &= \pi_j.
    \end{align*}
\end{proof}

\begin{theorem}[Limit Distribution is Invariant on Finite State Spaces]
    If \(X\) is a Markov chain on a finite state space \(I\), and
    \[
        \Prob\pqty{X_m = i} \to \pi_i
    \]
    as \(n \to \infty\) for \(i \in I\), then \(\pi\) is an invariant distribution for \(X\).
\end{theorem}

\begin{proof}
    Note that
    \[
        \Prob\pqty{X_{n + 1} = j} = \sum_{i \in I} \Prob\pqty{X_n = i} P_{ij}.
    \]

    Since the sum is finite, we can take \(n \to \infty\), giving
    \[
        \pi_j = \sum_{i \in I} \pi_i P_{ij},
    \]
    so \(\pi\) is invariant.

    Also, since \(\sum_{i \in I} \Prob\pqty{X_n = i} = 1\) for all \(n\), taking \(n \to \infty\) inside the sum, we have \(\sum_{i \in I} \pi_i = 1\), so \(\pi\) is a probability distribution.
\end{proof}

\begin{remarks}
    \begin{itemize}
        \item Here, the assumption that \(I\) is finite is crucial, since for example, the SSRW on \(\ZZ\), we have
        \[
            p_{xy} \pqty{n} \to 0
        \]
        for all \(x, y\) as \(n \to \infty\).

        \item If \(I\) is finite, then \(X\) always has an invariant distribution, by compatness. (We are not going to prove this.) Otherwise, things are more complex.
    \end{itemize}
\end{remarks}

\begin{example}
    We not recall the binary Markov chain we saw previously, with transition matrix
    \[
        P = \pmqty{
            1 - \alpha & \alpha \\
            \beta & 1 - \beta
        }.
    \]

    Then note
    \[
        p_{00} \pqty{n} = \frac{\beta}{\alpha + \beta} + \frac{\alpha}{\alpha + \beta} (1 - \alpha - \beta)^n \to \frac{\beta}{\alpha + \beta},
    \]
    if \(\alpha + \beta > 0\). So we have a limiting distribution
    \[
        \pqty{\frac{\beta}{\alpha + \beta}, \frac{\alpha}{\alpha + \beta}},
    \]
    and it is indeed invariant under the transition matrix.
\end{example}

\begin{definition}[Expected Number of Visits Before Return]
    For a fixed \(k \in I\), define
    \[
        \nu_k \pqty{j} = \Expt_k \pqty{\sum_{n = 0}^{T_k - 1} \indicator\pqty{X_n = j}}
    \]
    where \(T_k = \inf\set{n \geq 1}{X_n = k}\) is the first return time to \(k\).
\end{definition}

\begin{theorem}[Properties of \(\nu_k\)]
    Suppose the transition matrix \(P\) is irreducible and recurrent. Then we have the following.

    \begin{enumerate}[label=(\roman*)]
        \item \(\nu_k \pqty{k} = 1\) for all \(k \in I\).
        \item \(0 < \nu_k \pqty{i} < \infty\) for all \(i, k \in I\).
        \item \(\nu_k\) is an invariant measure for \(P\).
    \end{enumerate}
\end{theorem}

\begin{proof}
    \begin{enumerate}[label=(\roman*)]
        \item This is obvious from the definition, since we are counting the number of visits to \(k\) before returning to \(k\), which is exactly 1.
        \item[(iii)] Since \(P\) is recurrent, \(T_k < \infty\) with probability 1, and \(X_{T_k} = k\), we have
        \begin{align*}
            \nu_k \pqty{i} &= \Expt_k \pqty{\sum_{n = 0}^{T_k - 1} \indicator\pqty{X_n = j}} \\
            &= \Expt_k \pqty{\sum_{n = 1}^{T_k} \indicator\pqty{X_n = j}} \\
            &= \Expt_k \pqty{\sum_{n = 1}^{\infty} \indicator\pqty{X_n = i, n \leq T_k}} \\
            &= \sum_{n \geq 1} \Prob_k \pqty{X_n = i, n \leq T_k} \\
            &= \sum_{n \geq 1} \sum_{j \in I} \Prob_k \pqty{X_n = i, X_{n - 1} = j, n \leq T_k} \\
            &= \sum_{n \geq 1} \sum_{j \in I} \Prob_k \pqty{\Conditional{X_n = i}{X_{n - 1} = j, n \leq T_k}} \Prob_k \pqty{X_{n - 1} = j, n \leq T_k}.
        \end{align*}

        Note that
        \begin{align*}
            \Bqty{T_k \geq k} &= \Bqty{T_k \leq n - 1}\compl \\
            &= \Bqty{\bigcup_{t = 1}^{n - 1} \Bqty{X_k = t}}\compl,
        \end{align*}
        which only depends on \(X_0, X_1, \ldots, X_{n - 1}\), so it is a stopping time. Hence, by the strong Markov property,
        \begin{align*}
            \nu_k \pqty{i} &= \sum_{n \geq 1} \sum_{j \in I} P\pqty{j, i} \Prob_k \pqty{X_{n - 1} = j, T_k \geq n} \\
            &= \sum_{j \in I} P\pqty{j, i} \sum_{n \geq 1} \Expt_k \pqty{\indicator\pqty{X_{n - 1} = j, T_k \geq n}} \\
            &= \sum_{j \in I} P\pqty{j, i} \Expt_k \pqty{\sum_{n \geq 0} \indicator\pqty{X_{n} = j, T_k \geq n + 1}} \\
            &= \sum_{j \in I} P\pqty{j, i} \Expt_k \pqty{\sum_{n = 0}^{T_k - 1} \indicator\pqty{X_n = j}} \\
            &= \sum_{j \in I} P\pqty{j, i} \nu_k \pqty{j}.
        \end{align*}
        so \(\nu_k\) is an invariant measure.
        \item[(ii)] Suppose (iii) is true for this proof. Then we have \(\nu_k = \nu_k P\), and hence \(\nu_k = \nu_k P^n\) for all \(n \geq 1\). Then, we have
        \begin{align*}
            \nu_k \pqty{i} &= \pqty{\nu_k P^n}_i \\
            &= \sum_{j \in I} \nu_k \pqty{j} p_{ji} \pqty{n},
        \end{align*}
        and by irreducibility, we can choose \(n\) such that \(p_{ki} \pqty{n} > 0\). Then,
        \[
            \nu_k \pqty{i} \geq \nu_k \pqty{k} p_{ki} \pqty{n} = p_{ki} \pqty{n} > 0.
        \]

        Also, for any \(m \geq 1\), by invariance, we have
        \begin{align*}
            1 &= \nu_k \pqty{k} \\
            &= \sum_{j \in I} \nu_k \pqty{j} p_{jk} \pqty{m} \\
            &\geq \nu_k \pqty{i} p_{ik} \pqty{m},
        \end{align*}
        where we choose \(m \geq 1\) such that \(p_{ik} \pqty{m} > 0\) by irreducibility. Hence,
        \[
            \nu_k \pqty{i} \leq \frac{1}{p_{ik} \pqty{m}} < \plusinfty.
        \]
    \end{enumerate}
\end{proof}

\begin{theorem}[Uniqueness of Invariant Measure]
    \label{thm:uniqueness-of-invariant-measure}%
    Suppose \(P\) is irreducible and \(\lambda\) is an invariant measure such that \(\lambda_k = 1\). Then \(\lambda \geq \nu_k\) in components. Moreover, if \(P\) is recurrent, then \(\lambda = \nu_k\).
\end{theorem}

\begin{proof}
    Since \(\lambda\) is invariant, we have
    \begin{align*}
        \lambda_i &= \sum_{j \in I} \lambda_j P\pqty{j, i} \\
        &= P\pqty{k, i} + \sum_{j \in I \setminus \Bqty{k}} \lambda_j P\pqty{j, i} \\
        &= P\pqty{k, i} + \sum_{j_1 \neq k} \sum_{j_2 \in I} \lambda_{j_2} P\pqty{j_2, j_1} P\pqty{j_1, i} \\
        &= P\pqty{k, i} + \sum_{j_1 \neq k} P\pqty{k, j_1} P\pqty{j_1, i} + \sum_{j_1, j_2 \neq k} \lambda_{j_2} P\pqty{j_2, j_1} P\pqty{j_1, i} \\
        &= \cdots \\
        &= P\pqty{k, i} + \sum_{j_1 \neq k} P\pqty{k, j_1} P\pqty{j_1, i} + \cdots \\
        &\phantom{=} \qquad + \sum_{j_1, \ldots, j_{n - 1} \neq k} P\pqty{k, j_1} P\pqty{j_1, j_2} \cdots P\pqty{j_{n - 1}, i} \\
        &\phantom{=} \qquad + \sum_{j_1, \ldots, j_n \neq k} \lambda_{j_n} P\pqty{j_n, j_{n - 1}} \cdots P\pqty{j_1, i},
    \end{align*}
    and hence
    \[
        \lambda_i \geq \Prob_k \pqty{X_1 = i, T_k \geq 1} + \Prob_k \pqty{X_2 = i, T_k \geq 2} + \cdots + \Prob_k \pqty{X_n = i, T_k \geq n}.
    \]

    Therefore, we have
    \begin{align*}
        \lambda_i &\geq \sum_{n = 1}^{\infty} \Prob_k \pqty{X_n = i, T_k \geq n} \\
        &= \sum{n = 1}^{\infty} \Expt_k \pqty{\indicator\pqty{X_n = i, T_k \geq n}} \\
        &= \Expt_k \pqty{\sum_{n = 1}^{T_k} \indicator\pqty{X_n = i}} \\
        &= \nu_k \pqty{i}.
    \end{align*}

    If \(P\) is also recurrent, \(\nu_k\) is invariant, and hence so is the measure \(\lambda - \nu_k\), so
    \[
        \pqty{\lambda - \nu_k} = \pqty{\lambda - \nu_k} P^n
    \]
    for all \(n \geq 1\). The, for all \(i \in I\) and \(i \neq k\), we have
    \begin{align*}
        0 &= \pqty{\lambda - \nu_k}_k \\
        &= \bqty{\pqty{\lambda - \nu_k} P^n}_k \\
        &= \sum_{j \in I} \pqty{\lambda - \nu_k}_j p\pqty{j, k} \pqty{n}.
    \end{align*}

    Choose \(n \geq 1\) such that \(P_{ik} > 0\). Then
    \[
        0 \geq \pqty{\lambda - \nu_k}_i P_{ik} \pqty{n}
    \]
    and hence \(\pqty{\lambda - \nu_k}_i = 0\), so \(\lambda = \nu_k\).
\end{proof}

\begin{definition}[Positive and Null Recurrence]
    A recurrent state \(i\) is is \term{positive recurrent} if \(\Expt_i \pqty{T_i} < \infty\) , and \term{null recurrent} if \(\Expt_i \pqty{T_i} = \infty\).
\end{definition}

\begin{theorem}
    If \(P\) is irreducible, then the following are equivalent.

    \begin{enumerate}[label=(\roman*)]
        \item All states are positive recurrent.
        \item At least one state is positive recurrent.
        \item \(P\) has an invariant distribution \(\pi\).
    \end{enumerate}

    If any of these hold, then
    \[
        \pi_i = \frac{1}{\Expt_i \pqty{T_i}}
    \]
    for all \(i \in I\).
\end{theorem}

\begin{proof}
    \begin{itemize}
        \item (i) \(\implies\) (ii) is obvious.
        \item (ii) \(\implies\) (iii): Suppose \(k\) is positive recurrent. Then
        \begin{align*}
            \sum_{i \in I} \nu_k \pqty{i} = \sum_{i \in I} \Expt_k \pqty{\sum_{n = 0}^{T_k - 1} \indicator\pqty{X_n = i}} \\
            &= \Expt_k \pqty{T_k} < \infty,
        \end{align*}
        so define \(\pi = \flatfrac{\nu_k}{\Expt_k \pqty{T_k}}\). Since \(\nu_k\) is invariant, \(\pi\) is a probability distribution and invariant.
        \item (iii) \(\implies\) (i): Let \(k\) be arbitrary. Then
        \[
            \pi_k = \sum_{j \in I} \pi_j p_{jk} \pqty{n}
        \]
        and choosing a state \(i\) such that \(\pi_i > 0\) and \(n\) such that \(p_{ik} \pqty{n} > 0\), we have
        \[
            \pi_k \geq \pi_i p_{ik} \pqty{n} > 0,
        \]
        so \(\pi_k > 0\) for all \(k\).

        Now, for a fixed \(k\), define a new invariant measure
        \[
            \lambda_i = \frac{\pi_i}{\pi_k}.
        \]
        for all \(i\), so that \(\lambda_k = 1\). Then we must have that \(\lambda \geq \nu_k\) in components by \autoref{thm:uniqueness-of-invariant-measure}, so
        \begin{align*}
            \Expt_k\pqty{T_k} &= \sum_{i \in I} \nu_k \pqty{i} \\
            &\leq \sum_{i \in I} \lambda_i \\
            &= \frac{1}{\pi_k} \sum_{i \in I} \pi_i \\
            &= \frac{1}{\pi_k} < \infty
        \end{align*}
        hence \(k\) is positive recurrent.
    \end{itemize}

    Finally, since the chain is positive recurrent, we must have \(\lambda = \nu_k\), so
    \begin{align*}
        \frac{1}{\pi_k} &= \sum_{i \in I} \frac{\pi_i}{\pi_k} \\
        &= \sum_{i \in I} \lambda_i \\
        &= \sum_{i \in I} \nu_k \pqty{i} \\
        &= \Expt_k \pqty{T_k}.
    \end{align*}
\end{proof}